\documentclass[11pt,professionalfont,aspectratio=43]{beamer}

% Assets
\usepackage[czech]{babel}				% Jazyk
\usepackage[a-2u]{pdfx}				    % Kopírování z pdfka
\usepackage{tikz}						% Schémata automatů
\usetikzlibrary{shadows, shadows.blur, shapes.geometric, arrows, positioning, overlay-beamer-styles, patterns, shadings, shapes.callouts}
% \usetikzlibrary{arrows}
\usepackage{pgfplots}
% \pgfplotsset{compat=1.15}
% \usepackage{mathrsfs}
\usepackage[export]{adjustbox}

\usepackage[T1]{fontenc}
\usefonttheme[onlymath]{serif}
% \usepackage{newtxmath}
\usepackage{pxfonts}
\usepackage{slantsc}
% \usepackage{varwidth}

\usepackage[linesnumbered,ruled,vlined]{algorithm2e}                % Pseudokód algoritmů
\usepackage[utf8]{inputenc}
% \usepackage{textcomp}
\usepackage{hyperref}
\usepackage{fancybox}

%%% FONTS
% \usefonttheme{serif}
% \usepackage{lmodern}
\usepackage{helvet}

\hypersetup{%
    pdfencoding=auto,
    pdfauthor={\insertauthor},
    pdftitle={\insertsubtitle}
}
\usepackage{float}
\usepackage{csquotes}					% české uvozovky
\usepackage{enumerate}					% enumerate environment
\usepackage{indentfirst}
\usepackage{tcolorbox}                  % Fancy boxíky
\usepackage{mathtools}
\usepackage{pifont}
\usepackage{cancel}
\usepackage{subcaption}
\usepackage{soul}
\usepackage{makecell}                   % Zalomení v buňce tabulky
\usepackage{xcolor}
\usepackage{graphicx}
\usepackage{tabularx}
\usepackage{amsmath}
\usepackage{amsthm}
\usepackage{colortbl}                   % Barvení buněk v tabulce
\usepackage{multicol}                   % Rozdělení obsahu do sloucpů

\usepackage{listings}                   % Úryvky z kódu
\definecolor{maroon}{rgb}{0.5,0,0}
\definecolor{forestgreen}{rgb}{0.13, 0.55, 0.13}

% \lstdefinelanguage{XML}
% {
%   basicstyle=\ttfamily,
%   morestring=[s]{"}{"},
%   morecomment=[s]{?}{?},
%   morecomment=[s]{!--}{--},
%   commentstyle=\color{forestgreen},
%   moredelim=[s][\color{black}]{>}{<},
%   moredelim=[s][\color{red}]{\ }{=},
%   stringstyle=\color{blue},
%   identifierstyle=\color{maroon}
% }
% Colors

\def\maincolR{0} % červená složka (0–255)
\def\maincolG{83} % zelená složka
\def\maincolB{120} % modrá složka
\definecolor{maincolor}{HTML}{8B00B5}

% \pgfmathtruncatemacro{\authorfooterbgR}{1*\maincolR}
% \pgfmathtruncatemacro{\authorfooterbgG}{1*\maincolG}
% \pgfmathtruncatemacro{\authorfooterbgB}{1*\maincolB}
% \definecolor{authorfooterbg}{RGB}{\authorfooterbgR,\authorfooterbgG,\authorfooterbgB}

% \pgfmathtruncatemacro{\titlefooterbgR}{0.722*\maincolR}   
% \pgfmathtruncatemacro{\titlefooterbgG}{0.974*\maincolG}   
% \pgfmathtruncatemacro{\titlefooterbgB}{0.739*\maincolB}
% \definecolor{titlefooterbg}{RGB}{\titlefooterbgR,\titlefooterbgG,\titlefooterbgB}

% Beamer theme
\colorlet{themecolor}{maincolor}
\colorlet{authorfooterbg}{maincolor}
\definecolor{titlefooterbg}{HTML}{9471B0}
\definecolor{titlefooterfg}{HTML}{000000}
\colorlet{titlecolorfg}{maincolor}
\definecolor{titleboxbg}{HTML}{E7E7E7}
\definecolor{datefooterbg}{HTML}{E3E3E3}
\colorlet{titleframeboxcolor}{maincolor}

% Background gradient
\definecolor{gradMiddle}{RGB}{212,248,255}
\definecolor{gradEnd}{RGB}{114,231,255}

% Block theme
\definecolor{blocktitlebg}{HTML}{F0DFA9}
\definecolor{blocktitlefg}{HTML}{957300}
\definecolor{blockbodybg}{HTML}{F5F0DF}

\definecolor{alertblockbodybg}{HTML}{FFE4E4}
\definecolor{alertblocktitlebg}{HTML}{ECA4A4}
\definecolor{alertblocktitlefg}{HTML}{B11313}

\definecolor{exampleblockbodybg}{HTML}{DAF0DA}
\definecolor{exampleblocktitlebg}{HTML}{BBE0BB}
\definecolor{exampleblocktitlefg}{HTML}{257A25}

\definecolor{mathfg}{HTML}{00185E}

\colorlet{bulletbg}{maincolor}
% \definecolor{bulletbg}{HTML}{}

% Beamer theme
\usetheme{Madrid}
\usecolortheme{seahorse}
% \useoutertheme[subsection=false]{miniframes}

\usecolortheme[named=themecolor]{structure}
\setbeamertemplate{frame numbering}[fraction]
\setbeamertemplate{navigation symbols}{}
\setbeamercovered{invisible}

% Header/footer
\setbeamercolor{author in head/foot}{bg=authorfooterbg,fg=white}
\setbeamercolor{title in head/foot}{bg=titlefooterbg,fg=titlefooterfg}
\setbeamercolor{date in head/foot}{bg=datefooterbg,fg=black}
% \setbeamercolor{title}{bg=datefooterbg}
% \setbeamercolor{title}{fg=titlecolorfg}
\setbeamerfont{title}{series=\bfseries}
\setbeamercolor{frametitle}{bg=datefooterbg}
% \setbeamercolor{section in head/foot}{bg=red,fg=cyan}
% \setbeamercolor{subsection in head/foot}{bg=blue,fg=yellow}

% Table of contents
\setbeamercolor{section in toc}{fg=black}
\setbeamercolor{subsection in toc}{fg=red}
\setbeamercolor{section number projected}{fg=black}

% Blocks
\setbeamertemplate{blocks}[rounded][shadow=true]

\setbeamercolor{block title}{bg=blocktitlebg, fg=blocktitlefg}
\setbeamercolor{block body}{parent=normal text,use=block title,bg=blockbodybg}

\setbeamercolor{block body alerted}{bg=alertblockbodybg}
\setbeamercolor{block title alerted}{bg=alertblocktitlebg,fg=alertblocktitlefg}

\setbeamercolor{block body example}{bg=exampleblockbodybg}
\setbeamercolor{block title example}{bg=exampleblocktitlebg,fg=exampleblocktitlefg}

% Math
\setbeamercolor{math text}{fg=mathfg}

% Itemize a enumerate
\setbeamertemplate{itemize items}[circle]
\setbeamertemplate{enumerate items}[circle]
\setbeamercolor{itemize item}{fg=bulletbg,bg=white}
\setbeamercolor{itemize subitem}{fg=bulletbg,bg=white}
\setbeamercolor{itemize subsubitem}{fg=bulletbg,bg=white}

\setbeamertemplate{enumerate items}[circle]
\setbeamercolor{item projected}{bg=bulletbg}

% Titlepage
\defbeamertemplate*{title page}{customized}[1][]
{
    \vbox{}%
    \vfill%
    \begin{centering}
        \begin{tikzpicture}
            \node[thick, rectangle, rounded corners=5mm, drop shadow={shadow xshift=0mm, shadow yshift=2mm, opacity=0.2, blur shadow, shadow blur steps=10}, inner sep=5mm, fill=titleboxbg, fill opacity=1] {%
                {\usebeamerfont{title}\textcolor{titlecolorfg}{\inserttitle}\par}%
            };
        \end{tikzpicture}
        \vskip1em\par%
        \ifx\insertsubtitle\@empty%
        \else%
            {\usebeamerfont{subtitle}\insertsubtitle\par}%
        \fi%
        \vskip1em\par%
        \ifx\insertauthor\@empty%
        \else%
            {\usebeamerfont{author}\insertauthor}\\%
            \texttt{\email}
        \fi
        \vskip0.3em\par%
        % \ifx\insertinstitute\@empty%
        % \else%
        %     {\usebeamerfont{institute}\insertinstitute\par}%
        % \fi%
        \vskip1em\par%
        \ifx\inserttitlegraphic\@empty%
        \else%
            \vskip1em\par%
            \inserttitlegraphic\par%
        \fi%
        \vskip1em\par%
        \ifx\insertdate\@empty%
        \else%
            {\usebeamerfont{date}\insertdate\par}%
        \fi%
    \end{centering}%
    \vfill%
}
% Enumerate
%\setlist[enumerate]{topsep=0pt,itemsep=-1ex,partopsep=1ex,parsep=1ex,label=(\arabic*)}

\MakeOuterQuote{"}

%%% Makra pro definice, věty, tvrzení, příklady, ... (vyžaduje baliček amsthm)

\theoremstyle{plain}
\newtheorem{thm}{Věta}[section]
\newtheorem{lem}[theorem]{Lemma}
\newtheorem{prop}[theorem]{Tvrzení}
\newtheorem{cor}[theorem]{Důsledek}
\newtheorem*{prop*}{Tvrzení}

\theoremstyle{definition}
\newtheorem{define}[theorem]{Definice}
% \newtheorem{example}[theorem]{Příklad}
% \newtheorem{remark}[theorem]{Poznámka}
% \newtheorem{convention}[theorem]{Úmluva}
% \newtheorem{denoting}[theorem]{Značení}

\newenvironment{defblock}[1][]{
  \begin{block}{\ifthenelse{\equal{#1}{}}
  {Definice}
  {Definice (#1)}}
}{\end{block}}
\newenvironment{thmblock}[1][]{
  \begin{block}{\ifthenelse{\equal{#1}{}}
  {Věta}
  {Věta (#1)}}
}{\end{block}}
\newenvironment{corblock}[1][]{
  \begin{block}{\ifthenelse{\equal{#1}{}}
  {Důsledek}
  {Důsledek (#1)}}
}{\end{block}}
\newenvironment{propblock}[1][]{
  \begin{block}{\ifthenelse{\equal{#1}{}}
  {Tvrzení}
  {Tvrzení (#1)}}
}{\end{block}}
\newenvironment{exblock}[1][]{
  \begin{exampleblock}{\ifthenelse{\equal{#1}{}}
  {Příklad}
  {Příklad (#1)}}
}{\end{exampleblock}}


% Colors
\definecolor{lightblue}{HTML}{009AD4}
\definecolor{lightgreen}{HTML}{68FF00}
\definecolor{darkgreen}{HTML}{00B600}
\definecolor{darkblue}{HTML}{0265b3}
\definecolor{darkred}{HTML}{DA0707}
\definecolor{lightred}{HTML}{FF5100}
\definecolor{orange}{HTML}{FFE000}
\definecolor{codeblue}{HTML}{007eed}
\definecolor{codegreen}{rgb}{0,0.6,0}
\definecolor{codegray}{rgb}{0.5,0.5,0.5}
\definecolor{codebeige}{HTML}{D4A000}
\definecolor{codepink}{HTML}{FF0055}
\definecolor{backcolour}{rgb}{0.95,0.95,0.92}

\definecolor{decproblemtitlefg}{HTML}{957300}
\definecolor{decproblembg}{HTML}{F5F0DF}

\definecolor{timportantboxcolor}{HTML}{A97100}
\definecolor{talertboxcolor}{HTML}{B70000}
\definecolor{tnoticeboxcolor}{HTML}{159C00}
\definecolor{tinfoboxcolor}{HTML}{004DFF}

\newcommand{\markred}[1]{\textcolor{lightred}{#1}}
\newcommand{\markdarkred}[1]{\textcolor{darkred}{#1}}
\newcommand{\markgreen}[1]{\textcolor{lightgreen}{#1}}
\newcommand{\markdarkgreen}[1]{\textcolor{darkgreen}{#1}}
\newcommand{\markorange}[1]{\textcolor{orange}{#1}}
\newcommand{\markblue}[1]{\textcolor{lightblue}{#1}}
\newcommand{\markdarkblue}[1]{\textcolor{darkblue}{#1}}

% Math mode settings
% \AtBeginDocument{
%     \setlength{\abovedisplayskip}{3pt}
%     \setlength{\belowdisplayskip}{-4pt}
% }

% Inline images
\newcommand{\inlineimgscale}{1.1}

% X and check mark
\newcommand{\cmark}{\markgreen{\ding{51}}}
\newcommand{\xmark}{\markred{\ding{55}}}

% Math
\newcommand{\R}{\mathbb{R}}
\newcommand{\C}{\mathbb{C}}
\newcommand{\N}{\mathbb{N}}
\newcommand{\Z}{\mathbb{Z}}
\newcommand{\Q}{\mathbb{Q}}

\newcommand{\T}[1]{#1^\top}                     % Tranpozice matice
\newcommand{\compl}[1]{#1^\complement}
\newcommand{\mapping}[3]{#1:#2 \rightarrow #3}
\newcommand{\mapto}[3]{#1:#2 \mapsto #3}
\newcommand{\set}[1]{\left\{#1\right\}}
\newcommand{\powset}{\mathcal{P}}
\newcommand{\code}[1]{\langle#1\rangle}
\DeclareMathOperator*{\argmin}{arg\,min}
\DeclareMathOperator*{\argmax}{arg\,max}

% Statistika a AI
\newcommand{\average}[1]{\overline{#1}}
\DeclareMathOperator{\median}{med}
\DeclareMathOperator{\Var}{var}
\DeclareMathOperator{\Cov}{cov}
\DeclareMathOperator{\MSE}{MSE}
\DeclareMathOperator{\SSE}{SSE}

\DeclareMathOperator{\TP}{TP}
\DeclareMathOperator{\TN}{TN}
\DeclareMathOperator{\FP}{FP}
\DeclareMathOperator{\FN}{FN}
\DeclareMathOperator{\Accuracy}{Acc}
\DeclareMathOperator{\Precision}{Prec}
\DeclareMathOperator{\Recall}{Rec}
\DeclareMathOperator{\Fscore}{F}
\DeclareMathOperator{\Gini}{Gini}

\DeclareMathOperator{\mspan}{span}
\DeclareMathOperator{\mrank}{rank}
\DeclareMathOperator{\tg}{tg}
\DeclareMathOperator{\id}{id}
\DeclareMathOperator{\dom}{dom}
\DeclareMathOperator{\rng}{rng}
\renewcommand{\emptyset}{\varnothing}
\newcommand{\binary}[1]{\ensuremath{\left\langle #1\right\rangle}_B}

% DFS in and out lists
\DeclareMathOperator{\inlist}{in}
\DeclareMathOperator{\outlist}{out}
% \DeclareMathOperator{\lowlist}{low}
\DeclareMathOperator{\key}{key}

% A* macro
\newcommand{\Astar}{$\textcolor{black}{\textup{A}^{*}}$}

% Complexity classes
\DeclareMathOperator{\classP}{P}
\DeclareMathOperator{\classNP}{NP}
\DeclareMathOperator{\classTime}{TIME}
\DeclareMathOperator{\classNTime}{NTIME}
\DeclareMathOperator{\classSpace}{SPACE}
\DeclareMathOperator{\classNSpace}{NSPACE}

\DeclareMathOperator{\SAT}{SAT}
\DeclareMathOperator{\THREESAT}{3-SAT}
\DeclareMathOperator{\UNSAT}{UNSAT}
\DeclareMathOperator{\NZMNA}{NzMna}
\DeclareMathOperator{\KLIKA}{Klika}
\DeclareMathOperator{\ROZVRH}{Rozvrh}
\DeclareMathOperator{\SUDOKU}{Sudoku}
\newcommand{\HALT}{\ensuremath{\textsc{Halt}}}
\newcommand{\FACTOR}{\ensuremath{\textsc{Factor}}}

\DeclareMathOperator{\regex}{RegE}

%%% Barevné boxíky (notice / info / important / alert)
%
% Šířka boxu se nezadává ručně -- box se přizpůsobí šířce svého obsahu.
% Obsah se vysází do prostředí varwidth, takže krátký text zůstane na
% jednom řádku; delší text se zalomí, nejvýše však na šířku \linewidth
% (tj. \textwidth mimo sloupce a minipage). Pokud je titulek širší než
% obsah, řídí šířku boxu titulek.

\newsavebox{\tboxcontentbox}
\newsavebox{\tboxtitlebox}
\newlength{\tboxwidth}
\newlength{\tboxpadding}

% Geometrie boxu (musí odpovídat stylu tautobox níže):
% 2 * (boxrule + boxsep + left)
\setlength{\tboxpadding}{\dimexpr 0.5mm+1mm+4mm\relax}
\setlength{\tboxpadding}{2\tboxpadding}

\tcbset{tautobox/.style={boxrule=0.5mm, boxsep=1mm, left=4mm, right=4mm,
                         fonttitle=\bfseries}}

% Začátek sběru obsahu do boxu
\newcommand{\tboxcollect}{%
  \begin{lrbox}{\tboxcontentbox}%
    \begin{varwidth}{\dimexpr\linewidth-\tboxpadding\relax}%
}

% Ukončení sběru a výpočet výsledné šířky
% #1 = titulek (pro měření; prázdný u variant bez titulku)
\newcommand{\tboxmeasure}[1]{%
    \end{varwidth}%
  \end{lrbox}%
  \sbox{\tboxtitlebox}{\bfseries #1}%
  \setlength{\tboxwidth}{\wd\tboxcontentbox}%
  \ifdim\wd\tboxtitlebox>\tboxwidth
    \setlength{\tboxwidth}{\wd\tboxtitlebox}%
  \fi
  \addtolength{\tboxwidth}{\tboxpadding}%
  \ifdim\tboxwidth>\linewidth
    \setlength{\tboxwidth}{\linewidth}%
  \fi
}

% Vysázení boxu s titulkem: #1 = barva, #2 = titulek
\newcommand{\tboxrenderbox}[2]{%
  \tboxmeasure{#2}%
  \begin{center}%
    \begin{tcolorbox}[tautobox, colback=#1!15, colframe=#1,
                      title={#2}, width=\tboxwidth]
      \usebox{\tboxcontentbox}%
    \end{tcolorbox}%
  \end{center}%
}

% Vysázení boxu bez titulku: #1 = barva
\newcommand{\tboxrenderframe}[1]{%
  \tboxmeasure{}%
  \begin{center}%
    \begin{tcolorbox}[tautobox, colback=#1!15, colframe=#1,
                      width=\tboxwidth]
      \usebox{\tboxcontentbox}%
    \end{tcolorbox}%
  \end{center}%
}

% Poznámky (notice)
\newenvironment{tnoticebox}[1]{%
  \def\tboxtitletext{#1}\tboxcollect
}{%
  \tboxrenderbox{tnoticeboxcolor}{\tboxtitletext}%
}

\newenvironment{tnoticeframe}{%
  \tboxcollect
}{%
  \tboxrenderframe{tnoticeboxcolor}%
}

% Informační bloky (info)
\newenvironment{tinfobox}[1]{%
  \def\tboxtitletext{#1}\tboxcollect
}{%
  \tboxrenderbox{tinfoboxcolor}{\tboxtitletext}%
}

\newenvironment{tinfoframe}{%
  \tboxcollect
}{%
  \tboxrenderframe{tinfoboxcolor}%
}

% Důležité bloky (important)
\newenvironment{timportantbox}[1]{%
  \def\tboxtitletext{#1}\tboxcollect
}{%
  \tboxrenderbox{timportantboxcolor}{\tboxtitletext}%
}

\newenvironment{timportantframe}{%
  \tboxcollect
}{%
  \tboxrenderframe{timportantboxcolor}%
}

% Varovné bloky (alert)
\newenvironment{talertbox}[1]{%
  \def\tboxtitletext{#1}\tboxcollect
}{%
  \tboxrenderbox{talertboxcolor}{\tboxtitletext}%
}

\newenvironment{talertframe}{%
  \tboxcollect
}{%
  \tboxrenderframe{talertboxcolor}%
}

% Decision problem
\newcommand{\decproblem}[3]{
    \begin{center}
        \begin{minipage}{.9\textwidth}
            \begin{table}[!htbp]
                \centering
                \bgroup
                \def\arraystretch{1.8}
                \begin{tabularx}{\textwidth}{|rX|}
                \hline
                \rowcolor{decproblembg} \multicolumn{2}{|c|}{\textcolor{decproblemtitlefg}{\textsc{#1}}} \\ \hline
                \markdarkred{\textbf{Vstup:}} & #2                \\ \hline
                \rowcolor{decproblembg} \markdarkred{\textbf{Otázka:}} & #3               \\ \hline
                \end{tabularx}
                \egroup
            \end{table}
        \end{minipage}
    \end{center}
}
\newcommand{\bigO}{\mathcal{O}}

% TODO
\newcommand{\todo}[1]{\textcolor{red}{(\noindent TODO: #1.)}}


\newcommand{\hlcode}[1]{\colorbox{red}{#1}}

% Algorithms
\numberwithin{algocf}{section}
\renewcommand{\algorithmcfname}{Algoritmus}
\SetKwInput{KwIn}{Vstup}
\SetKwInput{KwOut}{Výstup}
\SetKwProg{Function}{Funkce}{}{end}
\setlength{\algomargin}{25pt}
\renewcommand{\thealgocf}{}     % Vypnutí číslování algoritmů

% Definice stylů pro uzly
\tikzstyle{startstop} = [rectangle, rounded corners, minimum width=1.5cm, minimum height=0.5cm,text centered, draw=black, fill=red!30]
\tikzstyle{process}   = [rectangle, minimum width=1.5cm, minimum height=0.5cm, text centered, draw=black, fill=orange!30]
\tikzstyle{decision}  = [diamond, aspect=2, text centered, draw=black, fill=green!30]
\tikzstyle{arrow}     = [thick,->,>=stealth]

\tikzstyle{vertex}    = [thick, draw, circle, minimum size=4mm, inner sep=1pt, outer sep=1pt, align=center]
\tikzstyle{vertexplain}    = [minimum size=4mm, inner sep=1pt, outer sep=1pt, align=center]
\tikzstyle{verteximg}    = [draw, circle, minimum size=4mm, inner sep=0pt, outer sep=1pt, align=center, clip]
\tikzstyle{verteximgplain}    = [minimum size=4mm, inner sep=0pt, outer sep=1pt, align=center]
\tikzstyle{task}    = [draw, rounded corners=2pt, minimum width=14mm, minimum height=8mm, inner sep=2pt, align=center]

\tikzstyle{edge}    = [thick,>=stealth]
\tikzstyle{diredge}    = [thick, ->, >=stealth]

%% Automaty
\tikzstyle{state}=[%
    draw,
    circle,
    thick,
    minimum size=10mm,
    inner sep=0pt,
    outer sep=0pt
]

\tikzstyle{acceptstate}=[%
    state,
    double,
    double distance=1pt
]

\tikzstyle{transition}=[%
    ->,
    >=stealth,
    thick
]

\newcommand<>{\state}[4][]{%
    \node#5[state,#1] (#2) at #3 {#4};
}

\newcommand<>{\acceptingstate}[4][]{%
    \node#5[acceptstate,#1] (#2) at #3 {#4};
}

\newcommand<>{\transitionarrow}[5][]{%
    \path#6 (#2) edge[transition,#1] node[#3] {#4} (#5);
}

\newcommand<>{\initialstate}[6][]{%
    \node#7[state,#1] (#2) at #3 {#4};
    \draw#7[transition] #5 -- (#2.#6);
}

%%%%%

% Frames
\newcommand{\tableofcontentsframe}{
    \begin{frame}{Obsah}
        \tableofcontents
    \end{frame}
}

% \definecolor{boxcolor}{HTML}{00820a}
\newcommand{\titleframe}[1]{%
  \begin{frame}[plain]
    \begin{tikzpicture}[remember picture, overlay]
      \node at (current page.center)
        [draw=titleframeboxcolor, line width=1pt,
         inner xsep=10mm, inner ysep=4mm, rounded corners=2mm,
         text width=0.8\textwidth, align=center,
         execute at begin node=\setlength{\baselineskip}{2.2em}]
        {\Huge\markred{#1}\normalfont};
    \end{tikzpicture}
  \end{frame}
  \addtocounter{framenumber}{-1}%
}

\newenvironment{titleframeenv}[1]{
    \begin{frame}[plain]
        \begin{tcolorbox}[colback=white,boxrule=1pt]
            \begin{center}
                \Huge\markred{#1}\normalfont
            \end{center}
        \end{tcolorbox}
}{\addtocounter{framenumber}{-1}\end{frame}}

\newcommand{\icon}[1]{
  \includegraphics[height=\fontcharht\font`\B]{#1}
}

% Obrázek, který se vždy vejde na snímek. Adjustbox jej v případě potřeby
% zmenší tak, aby nepřesáhl šířku textu ani zadanou část výšky snímku; díky
% tomu sedí sazba ve všech používaných poměrech stran (4:3, 16:9 i 16:10).
% Volitelným argumentem je maximální výška (výchozí 0.68\textheight), povinným
% cesta k souboru s obrázkem (TikZ zdroj vkládaný přes \input).
\newcommand{\obrazek}[2][0.68\textheight]{%
    \begin{figure}
        \centering
        \begin{adjustbox}{max width=\linewidth, max totalheight=#1, center}
            \input{#2}
        \end{adjustbox}
    \end{figure}%
}

% Totéž pro rastrové či vektorové obrázky vkládané přes \includegraphics.
\newcommand{\obrazekgrf}[2][0.68\textheight]{%
    \begin{figure}
        \centering
        \includegraphics[width=\linewidth, max totalheight=#1, center]{#2}
    \end{figure}%
}

% Závěrečný snímek se seznamem zdrojů. Prvním (volitelným) argumentem je název
% sekce i snímku, druhým čárkou oddělený seznam klíčů z ../assets/literatura.bib.
% Citace se sázejí podle ČSN ISO 690, v pořadí, v jakém jsou klíče uvedeny.
\newcommand{\zdrojeframe}[2][Zdroje]{%
    \section{#1}
    \nocite{#2}%
    \begin{frame}[allowframebreaks]{#1}
        \printbibliography[heading=none]
    \end{frame}
}
\input{../assets/listing.tex}

\def\presentationtitle{Aplikace nezávislé množiny a dalších problémů}
\def\presentationtitlefooter{Aplikace $\NZMNA$ a dalších}

\begin{document}

    \title[\presentationtitle]{\textsc{Teoretická informatika}}
\subtitle{\presentationtitle}
\author{David Weber}
\def\email{weber3@spsejecna.cz}
\institute{SPŠE Ječná}
\date{%
    Rok 
    \pgfmathsetmacro{\firstyear}{ifthenelse(\month<9, \year-1, \year)}%
    \pgfmathsetmacro{\secondyear}{ifthenelse(\month<9, \year, \year+1)}%
    \pgfmathtruncatemacro{\secondyearlasttwo}{mod(\secondyear,100)}%
    \pgfmathprintnumber[fixed,precision=0,zerofill,1000 sep={}]{\firstyear}/%
    \pgfmathprintnumber[fixed,precision=0,zerofill,1000 sep={}]{\secondyearlasttwo}%
}
\titlegraphic{\includegraphics[width=0.25\textwidth]{../images/SPSE-Jecna_Logo.pdf}}

\begin{frame}[plain]
    \titlepage
\end{frame}

    \section{Aplikace nezávislé množiny}

    \titleframe{Aplikace $\NZMNA$}

    \begin{frame}{Aplikace nezávislé množiny}
        \begin{itemize}
            \visible<2->{\item V mnoha praktických situacích máme k dispozici jediný sdílený zdroj a několik aktivit, které chtějí tento zdroj využívat v různých časových intervalech.}
            \visible<3->{\item Např. sestavování rozvrhu pro jednu učebnu, plánování schůzek jedné osoby nebo rezervace jednoho zařízení.}
            \visible<4->{\item Každá aktivita je určena začátkem a koncem svého trvání.}
            \visible<5->{\item Pokud se dvě aktivity časově překrývají, \textbf{nemohou být realizovány současně}.}
        \end{itemize}
    \end{frame}

    \begin{frame}{Rozvrhování}
        \begin{itemize}
            \visible<2->{\item Máme $n$ aktivit s přiděleným časovým intervalem.}
            \visible<3->{\item Časové intervaly budeme modelovat jako \textbf{intervaly reálných čísel}, tj. každá aktivita má interval
            \[I_j=(a_j,b_j)\; ,\; a_j,b_j\in\R,a_j<b_j,\]
            pro $j\in\set{1,\dots,n}$.}
        \end{itemize}
        \visible<4->{\begin{timportantframe}[0.8\textwidth]
            \begin{center}
                Které aktivity lze vybrat tak, aby si vzájemně časově nekolidovaly?
            \end{center}
        \end{timportantframe}}
    \end{frame}

    \begin{frame}{Formální pojetí problému}
        \begin{itemize}
            \visible<2->{\item Na vstupu máme množinu otevřených intervalů
            \[\mathcal{I}=\set{I_1,I_2,\dots,I_n},\]}
            \visible<3->{\item Pokud spolu dvě úlohy s intervaly $I_k$ a $I_\ell$ kolidují, pak $I_k\cap I_\ell\neq\emptyset$.}
            \visible<4->{\item Zde se nám hodí, že jsou intervaly otevřené, neboť pokud dvojice aktivit na sebe přímo navazuje, tj.
            \[I=(a,b)\quad\text{a}\quad J=(b,c),\]
            pak je jejich \textbf{průnik prázdný}, což dává smysl.}
        \end{itemize}
    \end{frame}

    \begin{frame}{Formální pojetí problému}
        \visible<2->{\decproblem{Rozvrhování}{
            Systém otevřených intervalů $\mathcal{I}=\set{I_1,I_2,\dots,I_n}$ a číslo $k\in\N$.
        }{
            Existuje podsystém $\mathcal{J}\subseteq\mathcal{I}$ takový, že $|\mathcal{J}|\geqslant k$ a pro každou dvojici intervalů $I,J\in\mathcal{J}$ platí $I\cap J=\emptyset$?
        }}
    \end{frame}

    \begin{frame}{Intervalový graf}
        \visible<2->{\begin{defblock}[Intervalový graf]
            Mějme systém otevřených intervalů $\mathcal{I}=\set{I_1,I_2,\dots,I_n}$. Pak \markdarkred{intervalový graf} $G_\mathcal{I}=(V,E)$, kde $V=\set{v_1,v_2,\dots,v_n}$ je graf splňující:
            \begin{itemize}
                \item vrchol $v_j$ odpovídá intervalu $I_j$,
                \item pro každé $k,\ell$ platí $v_kv_\ell\in E\iff I_k\cap I_\ell\neq\emptyset$.
            \end{itemize}
        \end{defblock}}
    \end{frame}

    \begin{frame}{Intervalový graf}{Příklad}
        \begin{figure}
            \centering
            \begin{tikzpicture}[x=0.7cm,y=0.7cm]

% -------------------------------------------------
% Intervaly vlevo
% -------------------------------------------------

\node at (4,5.9) {\small Intervaly};

% osa
\draw[arrow] (0,0) -- (8.3,0);
\foreach \x in {0,1,...,8}
{
    \draw (\x,0.12) -- (\x,-0.12);
    \node[below] at (\x,-0.12) {\small \x};
}

% interval I_1 = (0.5,2.5)
\node[left] at (0.5,1) {$I_1$};
\draw[darkblue, line width=1.4pt] (0.5,1) -- (2.5,1);
\draw[darkblue, line width=1pt, fill=white] (0.5,1) circle (2pt);
\draw[darkblue, line width=1pt, fill=white] (2.5,1) circle (2pt);

% interval I_2 = (1.5,4.2)
\node[left] at (1.5,2) {$I_2$};
\draw[darkblue, line width=1.4pt] (1.5,2) -- (4.2,2);
\draw[darkblue, line width=1pt, fill=white] (1.5,2) circle (2pt);
\draw[darkblue, line width=1pt, fill=white] (4.2,2) circle (2pt);

% interval I_3 = (3.2,5.4)
\node[left] at (3.2,3) {$I_3$};
\draw[darkblue, line width=1.4pt] (3.2,3) -- (5.4,3);
\draw[darkblue, line width=1pt, fill=white] (3.2,3) circle (2pt);
\draw[darkblue, line width=1pt, fill=white] (5.4,3) circle (2pt);

% interval I_4 = (5.0,7.0)
\node[left] at (5.0,4) {$I_4$};
\draw[darkblue, line width=1.4pt] (5.0,4) -- (7.0,4);
\draw[darkblue, line width=1pt, fill=white] (5.0,4) circle (2pt);
\draw[darkblue, line width=1pt, fill=white] (7.0,4) circle (2pt);

% interval I_5 = (2.2,6.3)
\node[left] at (2.2,5) {$I_5$};
\draw[darkred, line width=1.4pt] (2.2,5) -- (6.3,5);
\draw[darkred, line width=1pt, fill=white] (2.2,5) circle (2pt);
\draw[darkred, line width=1pt, fill=white] (6.3,5) circle (2pt);

% -------------------------------------------------
% Intervalový graf vpravo
% -------------------------------------------------

\node at (12.2,5.9) {\small Intervalový graf};

\node[draw=darkblue, circle, thick, minimum size=8mm, inner sep=0pt] (v1) at (10.0,4.1) {$I_1$};
\node[draw=darkblue, circle, thick, minimum size=8mm, inner sep=0pt] (v2) at (12.1,5.0) {$I_2$};
\node[draw=darkblue, circle, thick, minimum size=8mm, inner sep=0pt] (v3) at (14.1,4.1) {$I_3$};
\node[draw=darkblue, circle, thick, minimum size=8mm, inner sep=0pt] (v4) at (14.0,2.0) {$I_4$};
\node[draw=darkred, circle, thick, minimum size=8mm, inner sep=0pt] (v5) at (11.7,2.0) {$I_5$};

\draw[thick] (v1) -- (v2);
\draw[thick] (v1) -- (v5);
\draw[thick] (v2) -- (v3);
\draw[thick] (v2) -- (v5);
\draw[thick] (v3) -- (v4);
\draw[thick] (v3) -- (v5);
\draw[thick] (v4) -- (v5);

\end{tikzpicture}
        \end{figure}
    \end{frame}

    \begin{frame}{Intervalový graf}{Příklad: Největší množina nekolidujících aktivit}
        \begin{figure}
            \centering
            \begin{tikzpicture}[
    x=0.7cm,y=0.7cm,
    vertex/.style={
        draw=darkblue,
        thick,
        circle,
        minimum size=9mm,
        inner sep=1pt,
        font=\tiny,
        fill=white
    },
    selvertex/.style={
        draw=darkgreen,
        thick,
        circle,
        minimum size=9mm,
        inner sep=1pt,
        font=\tiny,
        fill=lightgreen!35
    },
    edge/.style={thick}
]

% -------------------------------------------------
% Zvýrazněná největší nezávislá množina
% -------------------------------------------------

\node[selvertex] (v1) at (0,0)   {$(0,1)$};
\node[selvertex] (v2) at (2.4,0) {$(1,2)$};
\node[selvertex] (v3) at (4.8,0) {$(2,3)$};
\node[selvertex] (v4) at (7.2,0) {$(3,4)$};
\node[selvertex] (v5) at (9.6,0) {$(4,5)$};

% -------------------------------------------------
% Ostatní vrcholy
% -------------------------------------------------

\node[vertex] (v6)  at (1.2,1.8) {$(0.6,1.4)$};
\node[vertex] (v7)  at (3.6,2.4) {$(1.6,2.4)$};
\node[vertex] (v8)  at (6.0,2.4) {$(2.6,3.4)$};
\node[vertex] (v9)  at (8.4,1.8) {$(3.6,4.4)$};
\node[vertex] (v10) at (4.8,4.6) {$(1.2,3.8)$};

% -------------------------------------------------
% Hrany bez křížení
% -------------------------------------------------

% "lokální" překryvy
\draw[edge] (v6) -- (v1);
\draw[edge] (v6) -- (v2);

\draw[edge] (v7) -- (v2);
\draw[edge] (v7) -- (v3);

\draw[edge] (v8) -- (v3);
\draw[edge] (v8) -- (v4);

\draw[edge] (v9) -- (v4);
\draw[edge] (v9) -- (v5);

% hrany z vrcholu v10 k dolním vrcholům
\draw[edge] (v10) to[out=215,in=90] (v2);
\draw[edge] (v10) -- (v3);
\draw[edge] (v10) to[out=325,in=90] (v4);

% hrany z vrcholu v10 k prostřední vrstvě
\draw[edge] (v10) to[out=210,in=85] (v6);
\draw[edge] (v10) to[out=235,in=95] (v7);
\draw[edge] (v10) to[out=305,in=85] (v8);
\draw[edge] (v10) to[out=330,in=95] (v9);

\end{tikzpicture}
        \end{figure}
        \begin{talertframe}[0.9\textwidth]
            \begin{center}
                Nekolidující aktivity tvoří v intervalovém grafu \textbf{nezávislou množinu}.
            \end{center}
        \end{talertframe}
    \end{frame}

    \begin{frame}{Převod $\ROZVRH\rightarrow\NZMNA$}
        \begin{itemize}
            \visible<2->{\item Na vstupu máme systém otevřených intervalů $\mathcal{I}$ a číslo $k\in\N$.}
            \visible<3->{\item Sestrojíme intervalový graf $G_\mathcal{I}$.}
            \visible<4->{\begin{timportantframe}[0.8\textwidth]
                \begin{center}
                    V grafu $G_\mathcal{I}$ existuje nezávislá množina velikosti alespoň $k$ právě tehdy, když lze vybrat alespoň $k$ vzájemně disjunktních intervalů z $\mathcal{I}$.
                \end{center}
            \end{timportantframe}}
            \visible<5->{\item Převod je zjevně polynomiální, neboť graf $G_\mathcal{I}$ má maximálně $\bigO(n^{2})$ hran.}
            \visible<6->{\item Celkově
            \[\ROZVRH(\mathcal{I},k)=\NZMNA(G_\mathcal{I},k).\]}
        \end{itemize}
    \end{frame}

    \begin{frame}{Převod $\NZMNA\rightarrow\ROZVRH$}
        \begin{itemize}
            \visible<2->{\item Převod $\NZMNA\rightarrow\ROZVRH$ nelze provést.}
            \visible<3->{\item Obecný graf $G$ totiž \textbf{nemusí nutně představovat} nějaký intervalový graf.}
            \visible<4->{\item Např. cyklus $C_4$ nemůže představovat intervalový graf.}
        \end{itemize}
    \end{frame}

    \begin{frame}{$C_4$ není intervalový graf}
        \visible<2->{\begin{figure}
            \centering
            \begin{tikzpicture}[x=0.7cm,y=0.7cm]

% -------------------------------------------------
% Vlevo: graf C_4
% -------------------------------------------------

\node at (2,5.6) {\small Graf $C_4$};

\node[draw=darkblue, circle, thick, minimum size=9mm, inner sep=0pt] (a) at (2,4) {$a$};
\node[draw=darkblue, circle, thick, minimum size=9mm, inner sep=0pt] (b) at (0,2) {$b$};
\node[draw=darkblue, circle, thick, minimum size=9mm, inner sep=0pt] (c) at (2,0) {$c$};
\node[draw=darkblue, circle, thick, minimum size=9mm, inner sep=0pt] (d) at (4,2) {$d$};

\draw[thick] (a) -- (b);
\draw[thick] (b) -- (c);
\draw[thick] (c) -- (d);
\draw[thick] (d) -- (a);

% chybějící hrany
% \draw[darkred, dashed, thick] (a) -- (c);
% \draw[darkred, dashed, thick] (b) -- (d);

% \node[darkred] at (1,2) {\small není hrana};
% \node[darkred] at (3,2) {\small není hrana};

% -------------------------------------------------
% Vpravo: pokus o intervalovou reprezentaci
% -------------------------------------------------

\node at (11,5.6) {\small Pokus o reprezentaci intervaly};

% osa
\draw[arrow] (6,0) -- (16,0);

% I_a
\node[left] at (8,1.2) {$I_a$};
\draw[darkblue, line width=1.4pt] (8,1.2) -- (14,1.2);
\draw[darkblue, line width=1pt, fill=white] (8,1.2) circle (2pt);
\draw[darkblue, line width=1pt, fill=white] (13,1.2) circle (2pt);

% I_b
\node[left] at (7,2.4) {$I_b$};
\draw[darkblue, line width=1.4pt] (7,2.4) -- (9.8,2.4);
\draw[darkblue, line width=1pt, fill=white] (7,2.4) circle (2pt);
\draw[darkblue, line width=1pt, fill=white] (9.8,2.4) circle (2pt);

% I_d
\node[left] at (11.3,3.6) {$I_d$};
\draw[darkblue, line width=1.4pt] (11.3,3.6) -- (14.2,3.6);
\draw[darkblue, line width=1pt, fill=white] (11.3,3.6) circle (2pt);
\draw[darkblue, line width=1pt, fill=white] (14.2,3.6) circle (2pt);

% I_c -- vede ke sporu
\node[left] at (9.2,4.8) {$I_c$};
\draw[darkred, line width=1.4pt] (9.2,4.8) -- (12.1,4.8);
\draw[darkred, line width=1pt, fill=white] (9.2,4.8) circle (2pt);
\draw[darkred, line width=1pt, fill=white] (12.1,4.8) circle (2pt);

% svislé pomocné čáry
\draw[densely dashed, gray] (9.2,0.2) -- (9.2,5.1);
\draw[densely dashed, gray] (9.8,0.2) -- (9.8,5.1);
\draw[densely dashed, gray] (11.3,0.2) -- (11.3,5.1);
\draw[densely dashed, gray] (12.1,0.2) -- (12.1,5.1);

\end{tikzpicture}
        \end{figure}}
        \visible<3->{\begin{talertframe}[0.8\textwidth]
            \begin{center}
                Interval $I_c$ má \textbf{neprázdný průnik} s $I_a$, ale v $C_4$ hrana $ac$ chybí.
            \end{center}
        \end{talertframe}}
    \end{frame}

\end{document}